\section{Representation by algebraic spaces}

\begin{definition}[Conjugate systems]
\label{akcp-7-1-conjugate}
\dcref{7.1}
\source{MR0555258-compactifying-picard:page-0048}
For locally finitely presented $\FF,\GG$ on $X/S$, with $\GG$ $S$-flat,
$\Conj(\FF,\GG)(T)$ consists of quotients of $\FF_T$ whose target is
$\GG_T\otimes\NN$ for an invertible $\NN$ on $T$.
\end{definition}

\begin{theorem}[Conjugate-system projective bundle]
\label{akcp-7-2-conjugate-representability}
\dcref{7.2}
\source{MR0555258-compactifying-picard:page-0048}
\uses{akcp-1-1-H,akcp-4-3-twist-descent,akcp-7-1-conjugate}
If $f:X\to S$ is proper finitely presented, $\GG$ is $S$-flat, and
$\OO_T\to(f_T)_*\mathcal H\!om(\GG_T,\GG_T)$ is an isomorphism for all $T$,
then $\Conj(\FF,\GG)$ is represented by an open retrocompact subscheme
$V\subset\PP(H(\FF,\GG))$.
\end{theorem}
\begin{proof}
As in \cref{akcp-4-2-linear-representability}, projective-bundle points are
maps $\FF_T\to\GG_T\otimes\NN$ nonzero on every fiber.  The conjugate quotient
condition is precisely that the cokernel be zero.  The zero-cokernel locus is
open and retrocompact by flattening, and the endomorphism hypothesis makes the
twist independent of choices by \cref{akcp-4-3-twist-descent}.
\end{proof}

\begin{lemma}[Regularity strata]
\label{akcp-7-3-regularity-strata}
\dcref{7.3}
\source{MR0555258-compactifying-picard:page-0049}
For a finitely presented projective $X/S$, a relatively very ample sheaf and an
$S$-flat finitely presented $\FF$, the locus $S_m$ where $\FF(s)$ is
$m$-regular is open retrocompact, $S_m\subset S_{m+1}$, and the $S_m$ cover
$S$.
\end{lemma}
\begin{proof}
Regularity is stable under increasing $m$, and Serre vanishing makes every
coherent fiber regular for some $m$.  After noetherian approximation, the
vanishing of the finitely many higher direct images characterizing
$m$-regularity is open by cohomology and base change; retrocompactness is then
automatic and descends.
\end{proof}

\begin{theorem}[Simple sheaves form an algebraic space]
\label{akcp-7-4-simple-space}
\dcref{7.4}
\source{MR0555258-compactifying-picard:page-0050}
\uses{akcp-2-7-local-quot,akcp-5-6-spl-sheaf,akcp-5-15-quot-strata-representable,akcp-7-2-conjugate-representability,akcp-7-3-regularity-strata}
For locally projective finitely presented $X/S$, the etale sheaf
$\Spl^{\acute{\mathrm et}}_{X/S}$ is represented by a quasi-separated
algebraic space locally of finite presentation over $S$.
\end{theorem}
\begin{proof}
The open subfunctors of simple $m$-regular sheaves with fixed polynomial cover
the sheaf by \cref{akcp-7-3-regularity-strata} and
\cref{akcp-5-8-spl-polynomial-open}.  On an affine base, every such sheaf is a
quotient of $E=\OO_X(-m)^{\oplus e}$; the corresponding open simple locus
$Z$ in the Quot scheme is representable by
\cref{akcp-5-15-quot-strata-representable}.  Conjugate quotients define a
Cartesian square with an open subscheme of a projective bundle by
\cref{akcp-7-2-conjugate-representability}; hence the equivalence relation on
$Z$ is smooth, finitely presented and effective.  Artin's quotient theorem
gives a quasi-separated algebraic space, and the quotient map is an etale
epimorphism.  The local quotients glue over $S$.
\end{proof}

\begin{remark}[Nonseparated simple-sheaf space]
\label{akcp-7-5-nonseparated}
\dcref{7.5}
\source{MR0555258-compactifying-picard:page-0051}
The simple-sheaf algebraic space need not be separated; families of simple
bundles of fixed rank and degree can have nonseparated specializations.
\end{remark}

\begin{corollary}[Picard algebraic space]
\label{akcp-7-6-picard-space}
\dcref{7.6}
\source{MR0555258-compactifying-picard:page-0051}
\uses{akcp-5-13-picard-open,akcp-7-4-simple-space}
If $\OO_X$ is $S$-simple, $\Pic^{\mathrm{CM},\acute{\mathrm et}}_{X/S}$ is
represented by an algebraic space locally finitely presented over $S$.
\end{corollary}
\begin{proof}
It is an open retrocompact subfunctor of the simple-sheaf sheaf by
\cref{akcp-5-13-picard-open}; take the corresponding open subspace in
\cref{akcp-7-4-simple-space}.
\end{proof}

\begin{remark}[A Picard functor need not be a scheme]
\label{akcp-7-7-picard-not-scheme}
\dcref{7.7}
\source{MR0555258-compactifying-picard:page-0051}
\uses{akcp-7-6-picard-space}
There are proper families for which the etale Picard functor is represented by
an algebraic space but by no scheme; the construction above still represents
it by the algebraic space of \cref{akcp-7-6-picard-space}.
\end{remark}

\begin{lemma}[Extending a generic rank-one sheaf]
\label{akcp-7-8-dvr-extension}
\dcref{7.8}
\source{MR0555258-compactifying-picard:page-0051}
\uses{akcp-3-3-embedding,akcp-3-4-hom-rank-one}
Let $S$ be the spectrum of a DVR with generic point $\eta$, and let
$X\to S$ be projective with integral fibers of the same dimension.  (i) Every
rank-one torsion-free sheaf on $X_\eta$ extends to a relatively rank-one
torsion-free sheaf on $X$.  (ii) Two such extensions whose generic fibers
become isomorphic after a field extension are already isomorphic.
\end{lemma}
\begin{proof}
Embed the generic sheaf in $\OO_{X_\eta}(m)$ by
\cref{akcp-3-3-embedding}.  Extend its cokernel flatly over $S$ and take the
kernel of $\OO_X(m)$ onto that extension.  Flatness of the family (regularity
of a DVR plus equal-dimensional integral fibers) makes the kernel $S$-flat,
and its closed fiber embeds in a line bundle, so it is torsion-free rank one.

For uniqueness, let $H=H(\II,\JJ)$.  The generic isomorphism gives a nonzero
section of $H_\eta$, hence a surjection $H\to\OO_S$ after splitting its free
and torsion parts.  The associated map $\II\to\JJ$ is nonzero on every fiber.
The two fibers have the same Hilbert polynomial, so
\cref{akcp-3-4-hom-rank-one} makes the map an isomorphism on each fiber; flatness
of $\JJ$ makes it an isomorphism globally.
\end{proof}

\begin{theorem}[Proper Picard algebraic space]
\label{akcp-7-9-picard-proper-space}
\dcref{7.9}
\source{MR0555258-compactifying-picard:page-0052}
\uses{akcp-3-4-hom-rank-one,akcp-5-13-picard-open,akcp-7-4-simple-space,akcp-7-8-dvr-extension}
For projective finitely presented $X/S$ with geometrically integral fibers of
fixed dimension, the torsion-free rank-one Picard etale sheaf is represented by
a proper finitely presented algebraic space over $S$; its pseudo-invertible
subsheaf is separated and finitely presented.
\end{theorem}
\begin{proof}
After noetherian approximation, the rank-one torsion-free fibers form a bounded
regular family by \cref{akcp-3-4-hom-rank-one}, so
\cref{akcp-7-4-simple-space} gives a quasi-separated locally finitely presented
space.  The bounded family makes it finitely presented.  The valuative
criterion is exactly \cref{akcp-7-8-dvr-extension}: existence gives universal
closedness and uniqueness gives separatedness.  Thus the torsion-free space is
proper; restricting to the open pseudo-invertible locus gives separatedness.
\end{proof}
